\section{Before you start}\label{sec:start}

\subsection{The reference: Stanford's Bump'em}
Bump'em was designed at Stanford. Everything they published is on one page:
\begin{center}
\url{https://biomechatronics.stanford.edu/bump-em}
\end{center}
It has the paper (Bump'em, ICRA 2020), the \textbf{build guide} (copy in
\file{docs/reference/stanford-bumpem-build-guide.pdf}), the \textbf{CAD files} for the printed motor parts,
the \textbf{bill of materials} with order numbers, the \textbf{ESCON configuration file}, and short
\textbf{assembly videos} for each step. This walkthrough follows the build guide and says, at each step,
where IBMS does it differently. Keep the guide open next to this document.

IBMS differences in one paragraph: force is measured with a DYMH-103 load cell and a Futek IAA100
amplifier (Stanford's ``closed-loop'' option, but with a bought load cell instead of the home-made strain-gauge
``dogbone''); one Teensy 4.1 with an MCP4728 DAC on a carrier board controls all four modules (Stanford
used a Speedgoat real-time computer); the towers are 90$\times$90 aluminium profiles; the software is our own
(firmware, Python server, browser UI). Sect.~\ref{sec:review} compares every part.

\subsection{The order of this walkthrough}
\begin{enumerate}
\item Build a tower: frame, motor unit, rope, force sensor, harness (Sect.~\ref{sec:walk-tower}).
\item Cabling: motor, force sensor, ground bridge, signal cable (Sect.~\ref{sec:walk-cabling}).
\item Power supplies: 48\,V, 24\,V, power strips, e-stop (Sect.~\ref{sec:walk-power}).
\item Chopper configuration, first power-on, ESCON setup (Sect.~\ref{sec:walk-chopper}).
\item Install the 3D-printed holders (Sect.~\ref{sec:walk-printed}).
\item Connect to the carrier PCB (Sect.~\ref{sec:walk-pcb}).
\item Download the software, run the server, first pulls (Sect.~\ref{sec:walk-software}).
\end{enumerate}
Do steps 1--4 for \emph{one} tower and get it working before building the next. Each step ends with a
\textbf{Check}; do not go on until it passes.

\subsection{What you need}
\begin{itemize}
\item \textbf{Tools} (Stanford guide p.\,9): metric and imperial hex keys, small flat and Phillips
screwdrivers, wire strippers and cutters, soldering iron, heat shrink. In addition: a \textbf{multimeter}
(continuity, DC volts, diode test), a ferrule crimper and a crimper for insulated terminals, a 4\,mm hex key
for the M6 button heads, a 3D printer (PETG or ASA).
\item \textbf{A Windows PC} for ESCON Studio, and any PC (Windows, macOS, Linux) for the Bump'em software.
\item \textbf{Parts}: Sect.~\ref{sec:bought} and \file{docs/parts-list.xlsx}.
\item \textbf{People}: the university electrician for everything on the 230\,V side.
\end{itemize}

\warn{Never work on wiring with a supply plugged in. The 230\,V terminals of the supplies (L, N, earth) are
electrician work only. The chopper holds its charge: before touching the 48\,V side, measure chopper pins
4--5 below 1\,V. First motor tests: nobody near the rope, hand on the power strip switch.}

\subsection{The version this walkthrough reaches}
\begin{table}[H]\centering\small
\begin{tabularx}{\linewidth}{lX}
\toprule
Part & Version and how to check it \\
\midrule
Firmware & v1.0.0, protocol 1 (\cmd{bumpem info} prints \cmd{proto=1,fw=1.0.0}) \\
Server & 0.4.0 (\cmd{GET /api/v1/info} prints \cmd{"server": "0.4.0"}) \\
Repositories & \fillin{git tag of both repos, e.g.\ \cmd{v2026-10-14}} \\
ESCON configuration & \file{escon/moduleA_left_2026-10-01.edc}, auto-tuned per driver \\
Carrier board & v1, \file{hardware/carrier/fab/bumpem-carrier-gerbers.zip} \\
Printed holders & \file{hardware/enclosures/out/*.stl}, \cmd{KEY\_CLEAR} 0.2\,mm \\
\bottomrule
\end{tabularx}
\end{table}
